\documentclass{article}

\usepackage[letterpaper]{geometry}
\usepackage[colorlinks=true, allcolors=blue]{hyperref}
\usepackage[spanish, activeacute]{babel}
\usepackage{tkz-euclide}
\usepackage{amsmath}
\usepackage{amsfonts}
\usepackage{amsthm}
\usepackage{amssymb,times}
\usepackage{graphicx}
\usepackage{latexsym}
\usepackage{subcaption}
\usepackage{xcolor}
\usepackage{pifont}
\usepackage{bbm}
\usepackage{pstricks}
\usepackage{pstricks-add}
\usepackage{pst-all}
\usepackage{mathrsfs}
\usepackage{pgfplots}
\usepackage{longtable}
\usepackage{multicol}
\usepackage{multirow}
\usepackage{kantlipsum}
\usepackage[inline]{enumitem}
\usepackage[export]{adjustbox}
\usepackage{array}
\usepackage{float}
\usepackage{booktabs}
\usepackage{mdframed}
\usepackage{minted}
\usepackage{tabularx}
\usepackage{adjustbox}
\usepackage[dvipsnames]{xcolor}

\usepackage{algorithm}
\usepackage{algpseudocodex}


\usepackage{caption}
\captionsetup[sub]{labelformat=empty}

\definecolor{bgcolor}{RGB}{240,245,255}

\usetikzlibrary{angles, quotes, positioning}
\renewcommand\thesection{\arabic{section}}
\renewcommand{\proofname}{Demostración}
\renewcommand{\qedsymbol}{$\blacksquare$}
\newcommand{\pa}[1]{\left(#1\right)}
\newtheorem{lemma}{Lema}

\newcommand{\R}{\mathbb{R}}
\newcommand{\N}{\mathbb{N}}
\newcommand{\Q}{\mathbb{Q}}
\newcommand{\Z}{\mathbb{Z}}

\newenvironment{solution}
{\par\noindent\textit{Solución.}\ }
{\hfill$\blacktriangle$\par}

\begin{document}
    
    \begin{table*}
        \centering
        \begin{tabular}{cc}
            \adjustbox{valign=c}{\includegraphics[width=0.12\linewidth, valign=m]{logo-ug.png}} & 
            \adjustbox{valign=c}{
                \begin{tabular}{c}
                    {\large\uppercase{\textbf{Universidad de Guanajuato}}}\\
                    {\small\uppercase{División de Ciencias Naturales y Exactas}}\\
                    {\small\uppercase{Campus Guanajuato}}
                \end{tabular}
            }
        \end{tabular}
    \end{table*}

    \begin{center}
        \textbf{Tarea 7 (Estructuras de Datos y Algoritmos I)}\\

        \begin{table*}
            \centering
            \begin{tabularx}{1\textwidth}{|X|X|X|}
                \hline
                \multicolumn{3}{|l|}{\textbf{Nombre:} Axel Magaña Falcón.}\\[2pt]
                \hline
                \textbf{Grupo: } 2$^\circ$ & \textbf{Fecha: }31/05/2026 & \textbf{Calificación:}\\[2pt]
                \hline
                \multicolumn{3}{|l|}{\textbf{Docente:} Claudia Esteves.}\\[2pt]
                \hline
            \end{tabularx}
        \end{table*}
    \end{center}

    \begin{enumerate}[label=\textbf{Problema \arabic*}]
        \setcounter{enumi}{3}
        \item \textbf{(Karamell)} \href{https://codeforces.com/gym/105327/submission/376832433}{\textcolor{ForestGreen}{\textbf{Accepted}} (Enlace a mi solución)}.
        
        Uma observación importante para resolver este problema es notar que si tenemos un par de conjuntos $A, B$ de bolsas de caramelos tales que 
        \begin{align*}
            \sum_{x \in A}x = \sum_{x \in B} x, 
        \end{align*}
        existe una forma de ordenar $A\cup B$ de forma que todas las bolsas de caramelos en $A$ las reciba Alice y todas las bolsas de caramelos en $B$ las reciba Bob, repartiéndolas según lo descrito en el problema. Véase el siguiente algoritmo en pseudocódigo que describe el proceso anterior:

        \begin{algorithm}[H]
            \begin{algorithmic}[1]
                \Procedure{OrderCaramels}{$A,B$}
                    \State $a \gets 0$
                    \State $b \gets 0$
                    \State $L \gets \varnothing$

                    \While{$A \neq \varnothing$ or $B \neq \varnothing$}
                        \If {$a \leq b$}
                            \State $a \gets a + A\text{\texttt{.back}}()$
                            \State $L\text{\texttt{.push\_back}}(A\text{\texttt{.back()}})$
                            \State $A\text{\texttt{.pop\_back}}()$
                        \Else
                            \State $b \gets b + B\text{\texttt{.back}}()$
                            \State $L\text{\texttt{.push\_back}}(B\text{\texttt{.back()}})$
                            \State $B\text{\texttt{.pop\_back}}()$
                        \EndIf
                    \EndWhile

                    \Return $L$ \Comment{$L$ será un ordenamiento válido}
                \EndProcedure
            \end{algorithmic}
        \end{algorithm}

        Así, nuestro problema se reduce a encontrar una partición del arreglo de bolsas de caramelos de dos subconjuntos cuyas sumas sean iguales. Llámese $A =\{a_1,a_2,\dots,a_n\}$ al conjunto de bolsas de caramelos. Sea $DP : \N\times\N \to 2^A, DP(i, k) \mapsto S$ donde $S \subset \{a_j : j \leq i\}$ es tal que $\sum_{x \in S} x = k$ para todo $1 \leq i \leq N$ y
        \begin{align*}
            1 \leq k \leq \frac{1}{2}\sum_{i=0}^N a_i
        \end{align*}
        Supongamos que $DP(i, k) = \varnothing$ si dicho subconjunto no existe. Al definir $DP$ como
        \begin{align*}
            DP(i, k) =
            \begin{cases}
                \{a_0\} &\text{si } i = 1,\,\,\text{y}\,\,a_0 = k,\\
                DP(i-1, k) &\text{si }i > 1\text{ y }DP(i-1,k) \neq \varnothing,\\
                DP(i-1, k-a_i) &\text{si }i > 1,\,\,k-a_i \geq 0\,\,\text{y}\,\,DP(i-1,k-a_i) \neq \varnothing,\\
                \varnothing &\text{si no se cumplen los casos anteriores,}
            \end{cases}
        \end{align*}
        podemos implementar la \texttt{DP} usando \texttt{bitset}. Con dicha solución, el problema se resuleve en ${\cal O}(N^2)$ (dado que $a_i$ es a lo más 100 y $N$ es a lo más 100).
    \end{enumerate}
\end{document}